% Secciones del juego: 5..11 (los cr\'editos son la 12, ver override-es.tex).

\sect{5. C\'omo jugar}

\textbf{La historia.} Algo maligno se esconde bajo la ciudad de Aguas Profundas (Waterdeep). Piergeiron, el primer Se\~nor de la ciudad, re\'une a un grupo de aventureros y le entrega una carta de encargo: bajar a las alcantarillas, descubrir qu\'e amenaza a la ciudad y destruirla. Al entrar, un derrumbe cierra el camino de vuelta. Desde ah\'i solo se puede ir hacia abajo: alcantarillas, ruinas enanas, las cuevas de los elfos oscuros y, al fondo, la guarida del contemplador Xanathar.

\textbf{El objetivo.} Recorrer los 12 niveles de la mazmorra, sobrevivir y acabar con Xanathar.

\textbf{Todo se ve desde los ojos del grupo.} La ventana grande muestra lo que tienes delante. Te mueves casilla a casilla y giras de 90 en 90 grados. Ver\'as pasar las paredes, objetos en el suelo, puertas, palancas, botones, inscripciones y monstruos que se acercan para luchar. Los monstruos no esperan: si te quedas quieto, ellos siguen movi\'endose y atacando.

\textbf{El grupo.} Empiezas con cuatro personajes. Los dos de arriba forman la \textbf{fila delantera}: son los \'unicos que pueden atacar cuerpo a cuerpo (espadas, mazas, hachas). Los dos de abajo van en la \textbf{fila trasera}: solo pueden atacar con armas a distancia (arcos, dagas o dardos lanzados) o con conjuros, pero siguen expuestos, especialmente a ataques laterales o por detrás. Por el camino encontrar\'as otros aventureros que pueden unirse al grupo, hasta un m\'aximo de seis personajes. Para \textbf{cambiar el orden}, haz clic derecho en el nombre de un personaje (en su lugar aparece \texttt{Cambiando}) y despu\'es clic derecho en el nombre de otro: los dos cambian de sitio. Un clic en cualquier otro punto del panel lo anula.

\textbf{Coger, usar, elegir.} La clave del juego es la diferencia entre estas tres acciones:
\begin{itemize}
\item \textbf{Coger} (bot\'on izquierdo): recoger o soltar objetos. Lo que coges pasa a ser el puntero del rat\'on y lo dejas donde quieras.
\item \textbf{Usar} (bot\'on derecho sobre una mano, o las teclas \key{1}--\key{6} y \key{Q}--\key{Y}): atacar con un arma, lanzar un conjuro con el libro o el s\'imbolo sagrado, beber una poci\'on, leer un pergamino\dots
\item \textbf{Elegir} (bot\'on izquierdo): pulsar los botones de los men\'us, elegir un conjuro de una lista, se\~nalar a un personaje.
\end{itemize}

\textbf{Examinar el mundo.} Haz clic con el bot\'on izquierdo sobre una inscripci\'on, un desag\"ue o un cad\'aver: en la zona de mensajes aparece lo que el grupo ve. Con el mismo clic se accionan palancas, botones y cadenas.

\textbf{Puertas.} No todas se abren igual. Muchas tienen un bot\'on o una palanca al lado (cl\'ic sobre \'el, o \key{ESPACIO} con la puerta delante). Las cerradas con llave se abren cogiendo la llave adecuada y haciendo clic con ella sobre la cerradura; un ladr\'on puede intentarlo con sus ganz\'uas del mismo modo. Una puerta a medio abrir se puede intentar forzar haciendo clic en su parte de abajo: lo intenta el personaje m\'as fuerte.

\textbf{Escaleras y portales} llevan a otros niveles. El juego recuerda lo que has dejado en cada nivel (puertas abiertas, monstruos muertos, objetos tirados) por si vuelves.

\sect{6. Crear el grupo}

Una partida nueva empieza en la pantalla \textbf{Generaci\'on de Caracteres}, donde creas a los cuatro aventureros. Todo se hace con el rat\'on, como en el original; con el teclado, los cursores eligen la ventana y \key{RETURN} la abre.

\begin{center}
\includegraphics[width=0.48\textwidth]{\shotdir chargen-race.png}\hfill
\includegraphics[width=0.48\textwidth]{\shotdir chargen-stats.png}
\end{center}

\begin{enumerate}
\item \textbf{Elige una ventana vac\'ia} de las cuatro de la izquierda.
\item \textbf{Raza y sexo, clase y alineamiento}, en este orden. Solo aparecen las clases que permite cada raza; \texttt{ATRAS} vuelve al paso anterior. Un grupo con un paladín no admite personajes malvados, y al rev\'es.
\item \textbf{Las caracter\'isticas} se tiran solas: fuerza, inteligencia, sabidur\'ia, destreza, constituci\'on y carisma, con la armadura (\texttt{AC}), los puntos de vida (\texttt{PV}) y el nivel. Antes, elige un \textbf{retrato}: se ven cuatro, y las flechas muestran otros.
\item Con los botones de la derecha: \texttt{REROLL} vuelve a tirar; \texttt{MODIFICA} deja subir o bajar cada valor con \texttt{+} y \texttt{-} dentro de los l\'imites de la raza y la clase (la que se cambia sale en rojo; con los cursores arriba y abajo pasas a otra; \texttt{OK} termina); \texttt{CARAS} cambia el retrato; \texttt{OK!} se queda con el personaje.
\item \textbf{Escribe su nombre} (hasta diez letras) y pulsa \key{RETURN}.
\end{enumerate}

Haz clic en una ventana con personaje para ver sus caracter\'isticas: \texttt{BORRAR} lo elimina para crear otro en su lugar. Con los cuatro creados aparece \texttt{JUGAR}: p\'ulsalo (o la tecla \key{P}) y cada personaje recibe el equipo de su clase (armadura, armas, el libro de conjuros del mago, el s\'imbolo sagrado del cl\'erigo, comida). El primero lleva adem\'as la carta de encargo de Aguas Profundas.

\subsection*{Elegir un grupo equilibrado}
Para empezar, combina dos combatientes delante con un clérigo y un mago detrás. El clérigo aporta curación y comida; el mago, ataques y protecciones. Un ladrón o una combinación de clases permite añadir ganzúas sin renunciar necesariamente a otra función. Lleva armas a distancia para la retaguardia.

\begin{reference}{Clase}{Papel en el grupo}
Guerrero & Especialista en armas y armadura; adecuado para recibir los golpes de frente.\\
Explorador & Combatiente que comparte muchas opciones de equipo del guerrero. No presupongas que dispone del libro de un mago.\\
Paladín & Combatiente bueno, con Imposición de manos y acceso a su selección de oraciones. Su alineamiento limita la composición del grupo.\\
Clérigo & Curaciones, alimento, protecciones y ayuda contra muertos vivientes. Puede llevar armadura, pero su elección de armas es más limitada.\\
Mago & Aprende conjuros de su libro y de pergaminos. Sus armas y protección física son limitadas: protégelo en la retaguardia.\\
Ladrón & Especialista en cerraduras con ganzúas. No todas las cerraduras pueden abrirse así; conserva también las llaves.\\
Varias clases & Un personaje reúne funciones de las clases indicadas. La experiencia se reparte entre ellas y su progreso no equivale al de una sola clase.\\
\end{reference}

\textbf{Razas.} Humanos, elfos, semielfos, enanos, gnomos y medianos ofrecen distintas combinaciones de clases y límites de características. El menú de creación muestra las posibilidades reales de cada elección. Un elfo no puede recibir el conjuro Resucitar de esta versión; tenlo presente antes de crear el grupo. El sexo y el retrato permiten definir al personaje, sin cambiar cómo se controla.

\begin{reference}{Característica}{Qué debes mirar}
Fuerza (STR) & Importante para el combate físico y para forzar puertas. Algunos guerreros muestran además una fuerza excepcional junto al 18.\\
Inteligencia (INT) & Característica del mago; interesa al elegir un personaje destinado a estudiar magia.\\
Sabiduría (WIS) & Característica del clérigo; interesa para su capacidad de preparar oraciones.\\
Destreza (DEX) & Ayuda a la protección del personaje. Comprueba su efecto en la CA, sobre todo si lleva poca armadura.\\
Constitución (CON) & Importante para la resistencia y los puntos de vida. No descuides los PV de quienes van delante.\\
Carisma (CHA) & Parte de la ficha y de los requisitos de algunas clases. No es una orden para negociar ni sustituye las respuestas de los diálogos.\\
\end{reference}

\textbf{Alineamiento.} Describe la orientación moral del personaje y restringe algunas clases y combinaciones del grupo. No añade un botón de habilidad. Puedes revisar los personajes antes de pulsar JUGAR; no hace falta aceptar la primera tirada.


\sect{7. La pantalla de aventura}

\begin{center}
\includegraphics[width=0.86\textwidth]{\imgdir hud.png}
\end{center}

\begin{tabularx}{\textwidth}{@{}clX@{}}
\toprule
\textbf{\tNo} & \textbf{\tElement} & \textbf{\tTellsYou}\\
\midrule
1 & Vista & Lo que el grupo tiene delante: paredes, puertas, objetos y monstruos. Clic izquierdo: coger, lanzar, dejar o examinar.\\
2 & Retrato y nombre & Qui\'en es cada personaje. Clic: su inventario. Gris: inconsciente; una calavera: muerto. Si recibe un golpe aparece una mancha roja con el da\~no.\\
3 & Mano de arriba & Lo que lleva en la mano principal, normalmente el arma. Clic derecho o \key{1}--\key{6}: usarla. Tras usarla queda en gris un momento.\\
4 & Barra de \texttt{PV} & Los puntos de vida: verde; amarilla cuando queda menos de un tercio; roja cuando llega a cero y el personaje cae.\\
5 & Mano de abajo & La mano secundaria: escudo, segunda arma, libro de conjuros, s\'imbolo sagrado\dots{} Clic derecho o \key{Q}--\key{Y}: usarla.\\
6 & Flechas & Avanzar, retroceder, girar y moverse de lado.\\
7 & Br\'ujula & Hacia d\'onde mira el grupo. Mientras se lanza un conjuro, aqu\'i aparece el men\'u de conjuros.\\
8 & Mensajes & Lo que ocurre: golpes, venenos, textos de la mazmorra, lo que dicen otros personajes.\\
9 & \texttt{CAMP} & El campamento: descansar, preparar conjuros, preferencias y guardar o cargar la partida.\\
\bottomrule
\end{tabularx}

\sect{8. El inventario y la hoja del personaje}

Haz clic en un retrato y el panel de la derecha muestra el inventario de ese personaje. Vuelve al juego haciendo clic otra vez en el retrato.

\begin{center}
\includegraphics[width=0.42\textwidth]{\imgdir inv.png}\hfill
\includegraphics[width=0.42\textwidth]{\shotdir stats-panel.png}
\end{center}

\begin{tabularx}{\textwidth}{@{}clX@{}}
\toprule
\textbf{\tNo} & \textbf{\tElement} & \textbf{\tTellsYou}\\
\midrule
1 & Retrato & Clic: vuelta a la pantalla de aventura.\\
2 & \texttt{HP} y \texttt{FOOD} & Puntos de vida y comida. Sin comida, el personaje pierde vida poco a poco y no puede preparar conjuros.\\
3 & Plato & Para comer: coge una raci\'on de comida y su\'eltala sobre el plato.\\
4 & Flechas & El personaje anterior o el siguiente (tambi\'en sirven para pasar un objeto de un personaje a otro).\\
5 & Mochila & Catorce casillas para llevar cosas.\\
6 & Carcaj & Las flechas. Suelta flechas sobre \'el para llenarlo; el n\'umero dice cu\'antas quedan.\\
7 & Mano principal & El arma (la misma mano de arriba del panel).\\
8 & Mano secundaria & Escudo, libro, s\'imbolo sagrado, munici\'on\dots\\
9 & Hoja & Pasa a la hoja del personaje (a la derecha): clase, alineamiento, raza, caracter\'isticas, experiencia y nivel.\\
10 & Cuerpo & Casillas para lo que se lleva puesto: cabeza, cuello, armadura, mu\~necas, anillos y botas; a la derecha, el cinto con tres casillas para objetos a mano (las armas arrojadizas que haya en \'el pasan solas a la mano al lanzar).\\
\bottomrule
\end{tabularx}

Para ponerte algo, c\'ogelo y su\'eltalo en su casilla. Si llevas un objeto en el puntero y haces clic en el retrato de un personaje, se abre su inventario para que lo guardes all\'i.

\subsection*{Leer la ficha y reconocer los estados}
\textbf{CA/AC} es la clase de armadura: cuanto menor, mejor, incluso si es negativa. \textbf{PV/HP} indica la vida actual y máxima. La experiencia permite subir de nivel; no es dinero ni una puntuación que se gaste. En personajes de varias clases se muestran sus niveles por separado. El nivel del personaje y el nivel de un conjuro son cosas distintas.

\begin{reference}{Estado}{Qué hacer}
Herido & Cura con magia, pociones o descanso seguro. La comida no repone directamente los PV.\\
Hambriento & Alimenta al personaje sobre el plato del inventario. El hambre impide recuperar conjuros y acaba costando vida.\\
Envenenado & Usa antídoto o Neutralizar veneno. Retrasar veneno solo da tiempo; dormir con veneno puede matar. Después cura el daño ya sufrido.\\
Paralizado & El personaje no puede actuar; algunos ataques lo dejan así unos dos minutos y medio. Protégelo y utiliza Quitar parálisis cuando esté disponible. Curar PV no elimina este estado.\\
Petrificado & No puede actuar. La curación no revierte la petrificación. Conserva un guardado anterior por si necesitas recuperar al grupo.\\
Inconsciente & Con PV de cero o menos, pero por encima de $-10$, el personaje ha caído. Una curación que lo lleve por encima de cero puede devolverlo al combate.\\
Muerto & A $-10$ PV aparece la calavera. Las curaciones ordinarias no sirven; Resucitar, cuando sea aplicable, lo devuelve con 1 PV. Si cae todo el grupo se intenta cargar la partida guardada.\\
\end{reference}

\textbf{El inventario no pausa la aventura.} Organiza el equipo en un lugar despejado. Para detenerte a pensar, abre el campamento; al descansar sí transcurre tiempo y vuelven a importar el hambre, los enemigos y el veneno.


\sect{9. Combate y magia}

\textbf{Atacar.} Usa (clic derecho, o \key{1}--\key{6}) el arma de la mano de un personaje de la fila delantera cuando el monstruo est\'e delante, en la vista. Las armas a distancia se usan igual desde cualquier fila: un arco necesita flechas en la otra mano, y al disparar se recargan solas del carcaj; las dagas y los dardos se reponen desde el cinto. Tras cada ataque la mano queda en gris hasta que el personaje est\'a listo otra vez. Al morir un monstruo, el grupo gana experiencia y los personajes suben de nivel.

\textbf{Lanzar conjuros.} Los magos usan su \textbf{libro de conjuros}; los cl\'erigos y paladines, su \textbf{s\'imbolo sagrado}. Usa el libro o el s\'imbolo (clic derecho) y en lugar de la br\'ujula aparece el men\'u de conjuros:

\begin{center}
\includegraphics[width=0.6\textwidth]{\shotdir spells.png}
\end{center}

Arriba est\'an los niveles de conjuro (\texttt{1RO} a \texttt{5TO}); debajo, los conjuros preparados de ese nivel. Haz clic en uno para lanzarlo; si afecta a un personaje, haz clic despu\'es en el retrato de ese personaje. Los conjuros de ataque solo se lanzan contra monstruos que est\'en a la vista. \texttt{ABORTAR} cierra el men\'u. Un conjuro lanzado se gasta: hay que volver a prepararlo descansando en el campamento.

\textbf{Pergaminos, pociones y varitas.} Se usan igual que un arma: en la mano y clic derecho. El pergamino se gasta al leerlo. Un mago puede copiar en su libro el conjuro de un pergamino (campamento, \texttt{Utilizar Pergaminos}).

\textbf{Expulsar muertos vivientes.} Los cl\'erigos, y los paladines a partir del nivel 3, intentan solos espantar a los no muertos que tengan delante.

\textbf{Venenos y par\'alisis.} Algunos monstruos envenenan o paralizan. El veneno quita vida poco a poco hasta que se cura (con un conjuro o una poci\'on); en la hoja del inventario aparece \texttt{ENVENENADO}.

\subsection*{Del libro a un conjuro listo: ejemplo completo}
\begin{enumerate}
\item Entra en CAMP y elige Memorizar Conjuros para un mago, u Orar Conjuros para un clérigo o paladín.
\item Selecciona el personaje y el nivel de conjuro. Añade copias con clic izquierdo o el botón +; quítalas con clic derecho o el botón $-$. Si indica 0 restantes, primero libera un hueco. Por ejemplo, cambia una Ayuda por un Paralizar.
\item Sal de la lista y elige Descansar grupo. Elegir conjuros no los recupera al instante. Espera a que los mensajes indiquen que han sido preparados; si interrumpes el descanso puede quedar trabajo pendiente.
\item Vuelve a la aventura, usa el libro o símbolo desde una mano, elige el nivel y el conjuro listo. Si pide un destinatario, señala su retrato.
\item Tras gastarlo, vuelve a descansar para recuperarlo. Para aprender una magia nueva de mago, primero copia el pergamino y después selecciónala para preparar.
\end{enumerate}

\textbf{Teclado en la preparación.} Las flechas arriba y abajo resaltan un conjuro; izquierda y derecha cambian de nivel. Enter o espacio añade una copia, la tecla menos quita una y Esc sale conservando la selección. El número de la derecha es la cantidad elegida, no un campo donde escribir cifras.

\textbf{Contacto, proyectil y protección.} Un proyectil necesita una dirección libre hacia el objetivo; un conjuro de contacto puede exigir primera fila o una mano libre. En Golpe eléctrico, Toque vampírico y Hoja llameante debes usar después la mano que contiene la magia. Una protección no hace daño por sí sola. La sección 13 explica cada caso.

\textbf{Ritmo del combate.} Ataca, espera a que la mano vuelva a estar disponible y vuelve a usarla. Retroceder o desplazarte de lado puede darte margen, siempre que hayas comprobado el suelo. No abras el inventario esperando detener a los enemigos. Recupera los proyectiles que hayan quedado en el suelo cuando la zona sea segura.


\sect{10. El campamento}

Haz clic en \texttt{CAMP}. El juego se detiene mientras est\'as en el campamento.

\begin{center}
\includegraphics[width=0.32\textwidth]{\shotdir camp.png}\hfill
\includegraphics[width=0.32\textwidth]{\shotdir memorize.png}\hfill
\includegraphics[width=0.32\textwidth]{\shotdir options.png}
\end{center}

\begin{itemize}
\item \texttt{Descansar grupo}: el grupo descansa hora a hora (la cuenta aparece en la ventana). Tras unas ocho horas cada mago y cada cl\'erigo empieza a recuperar sus conjuros, uno tras otro, y la ventana de mensajes lo cuenta. Si alguien est\'a herido y hay un cl\'erigo con conjuros de curaci\'on, el juego pregunta si quieres que cure al grupo. Cuando ya no quedan conjuros por recuperar pero alguien sigue herido, pregunta si seguir descansando hasta curarse. Cada ocho horas los personajes recuperan un punto de vida y comen. Un clic interrumpe el descanso. Con monstruos cerca no se puede descansar, y si se acerca alguno se sale del campamento con un aviso. Cuidado con el veneno: un personaje envenenado pierde vida cada hora y puede morir descansando (el juego avisa antes). Si el grupo tiene hambre, el juego pregunta si quieres seguir descansando.
\item \texttt{Memorizar Conjuros}: elige a un mago y escoge los conjuros que aprender\'a al descansar. Los n\'umeros de la derecha son cu\'antos tendr\'a de cada uno; \texttt{Borrar} quita la elecci\'on y los botones \texttt{1}--\texttt{5} cambian de nivel.
\item \texttt{Orar Conjuros}: lo mismo para cl\'erigos y paladines.
\item \texttt{Utilizar Pergaminos}: un mago copia en su libro un conjuro nuevo de un pergamino del grupo.
\item \texttt{Preferencias}: \texttt{Sonidos} activa o quita los efectos de sonido. La música suena durante la introducción.
\item \texttt{Opciones}: \texttt{Cargar Juego}, \texttt{Guardar Juego}, \texttt{Dejar Personaje} (retira un miembro cuando hay cinco o seis; detalles más abajo) y \texttt{Abandonar Juego} (vuelta a MSX-DOS).
\item \texttt{Salir}: vuelta al juego.
\end{itemize}

En los men\'us del campamento tambi\'en sirve el teclado: las flechas mueven el bot\'on iluminado, \texttt{RETURN} o la barra espaciadora lo eligen y \texttt{ESC} sale (en las preguntas S\'i/No las flechas izquierda y derecha cambian de respuesta). Cuando hay que escoger un personaje, las flechas mueven el marco de colores y \texttt{RETURN} lo elige. En la lista de conjuros, las flechas arriba y abajo (o el puntero) iluminan uno, \texttt{RETURN}, la barra espaciadora o un clic lo a\~naden, la tecla \texttt{-} quita uno y las flechas izquierda y derecha cambian de nivel.

\textbf{Guardar.} Hay una sola partida guardada para cada idioma: \texttt{Guardar Juego} sustituye a la anterior. Es la que se carga con \texttt{CARGAR ULTIMA PARTIDA} en la pantalla de t\'itulo y cuando todo el grupo cae.

\subsection*{Cada opción y sus botones}
\begin{reference}{Menú / opción}{Resultado}
Memorizar / Orar & Elige primero un personaje apto. Los botones 1--5 cambian el nivel del conjuro, no el nivel del personaje. Un clic izquierdo en una entrada añade una copia; un clic derecho la quita. Los botones + y $-$ actúan sobre la entrada resaltada. Con 0 huecos restantes, quita una copia antes de añadir otra.\\
Borrar & En preparación de conjuros vacía la selección de la página de nivel actual. No borra personajes ni elimina del libro los conjuros aprendidos.\\
Salir de la preparación & Conserva la selección y vuelve al campamento. Aún debes descansar para recuperar las copias pendientes.\\
Utilizar Pergaminos & Selecciona un mago y un conjuro nuevo que pueda copiar. Puede aprovechar pergaminos de mago llevados por los miembros del grupo. Copiar consume el pergamino y enseña el conjuro; no lo deja automáticamente preparado. Los pergaminos divinos no se copian.\\
Preferencias: Sonidos & Activa o desactiva los efectos si tienes una MoonSound conectada.\\
Preferencias: Salir & Vuelve al campamento.\\
Opciones: Cargar Juego & Recupera la partida guardada del idioma actual. Se pierde el progreso posterior a ese guardado. Si la carga falla, lee el mensaje antes de seguir.\\
Opciones: Guardar Juego & Escribe el estado actual y sustituye el guardado anterior del mismo idioma. Espera a que termine; no apagues durante la escritura.\\
Opciones: Dejar Personaje & Solo está disponible con cinco o seis miembros. Elige a quien sale: no puede ser el primero, pero no se limita a los compañeros reclutados. Sus pertenencias quedan en el suelo; revisa el grupo y recupera lo necesario.\\
Opciones: Abandonar Juego & Pide confirmación para volver a MSX-DOS. Sí abandona; No cancela. No guarda automáticamente.\\
Opciones: Salir & Regresa al campamento. El Salir del menú principal de CAMP reanuda la aventura.\\
\end{reference}

\textbf{Preguntas durante el descanso.} Aceptar curaciones permite utilizar las disponibles; seguir descansando para sanar prolonga el descanso y el consumo de comida. Aceptar descansar con hambre o veneno no elimina esos peligros. Un clic interrumpe el proceso; los monstruos pueden impedirlo o interrumpirlo también.

\textbf{Diálogos.} Lee el mensaje y elige la respuesta ofrecida: aceptar o rechazar ayuda, reclutar o responder a un personaje puede cambiar lo que ocurre. No todos los encuentros son combates. Si el grupo ya tiene seis miembros, necesitas una plaza para incorporar a otro.

\textbf{Qué se guarda.} Cada idioma tiene un guardado para el grupo y su progreso. Al elegir \texttt{Guardar Juego}, se sustituye el guardado anterior de ese idioma. Si creas un grupo nuevo, podrás recuperar el guardado anterior hasta que vuelvas a guardar.

\sect{11. Equipo y objetos de la aventura}

Los dibujos muestran un ejemplo de cada familia; las variantes pueden compartir icono. El color de una poción o un anillo no basta para identificar su efecto.

Cada objeto tiene un lugar de uso: mano, cuerpo, inventario o mecanismo de la mazmorra. El nombre no basta para saber sus propiedades. Al coger un pergamino con el botón izquierdo se muestra el nombre de su conjuro sin consumirlo. Las pociones y varitas sin identificar conservan nombres genéricos: Leer magia permite identificarlas.

\textbf{Bonificaciones y maldiciones.} Una bonificación positiva puede mejorar el ataque, el daño o la armadura. Una maldición puede empeorar el equipo y bloquear su retirada. Guarda antes de probar objetos sospechosos. No hay una opción de campamento para eliminar maldiciones. Las criaturas que oxidan o destruyen equipo pueden hacerte perder objetos: conserva recambios.

\refspace
\subsection*{Armas y munición}

Daño base contra enemigos pequeños/medianos y grandes, respectivamente. Se expresa antes de fuerza, bonificaciones y defensas. \textbf{G}: guerrero, explorador y paladín; \textbf{M}: mago; \textbf{C}: clérigo; \textbf{L}: ladrón. Las combinaciones de clase pueden ampliar las opciones; el juego comprueba si el personaje puede usar el equipo.

\begin{reference}{Nombre}{Efecto, uso y límites}
\itemname{axe}{Hacha} & 1d8 / 1d10. Una mano; G, L. Incluye Drow Cleaver.\\[4pt]
\itemname{long-sword}{Espada larga} & 1d8 / 1d12. Una mano; G, L. Variantes: Slasher, Night Stalker, Severious y espada de adamantita.\\[4pt]
\itemname{short-sword}{Espada corta} & 1d6 / 1d8. Una mano; G, L. Slicer pertenece a esta familia.\\[4pt]
\itemname{dagger}{Daga} & 1d4 / 1d3. G, M, L. Se lanza al usarla desde la mano; recógela después. Backstabber, Flicka y Guinsoo son variantes.\\[4pt]
\itemname{spear}{Lanza} & 1d6 / 1d8. G, L. Se usa como arma arrojadiza desde la mano. No se repone si no llevas otra.\\[4pt]
\itemname{halberd}{Alabarda} & 1d10 / 2d6. Dos manos; G, L. La alabarda del jefe es una variante de esta familia.\\[4pt]
\itemname{mace}{Maza / cetro} & 1d6 + 1 / 1d6. Una mano; G, C, L. El Cetro de poder real funciona aquí como un arma de esta familia; no presupongas un poder especial de empuje.\\[4pt]
\itemname{flail}{Mangual} & 1d6 + 1 / 2d4. Una mano; G, C, L. Para golpear desde la primera fila.\\[4pt]
\itemname{staff}{Bastón} & 1d6 / 1d6. Dos manos; todas las clases. También lo puede empuñar un mago.\\[4pt]
\itemname{bow}{Arco} & 1d6 / 1d6. Dos manos; G, L. Necesita flechas en la otra mano y puede recargarlas del carcaj. Incluye el arco drow.\\[4pt]
\itemname{sling}{Honda} & 1d4 / 1d4. Dos manos; G, C, L. Usa piedras como munición. Comprueba ambas manos si aparece el aviso de falta de munición.\\[4pt]
\itemname{dart}{Dardo} & 1d3 / 1d2. G, M, L. Arrojadizo; puede reponerse desde el cinto. Los dardos especiales tienen modificadores propios.\\[4pt]
\itemname{arrow}{Flecha y piedra} & Munición. Las flechas se guardan en el carcaj y alimentan el arco. Las piedras alimentan la honda y también pueden lanzarse. Son objetos recuperables cuando quedan en el suelo.\\[4pt]
\end{reference}

\refspace
\subsection*{Armadura y protección}

Una CA más baja es mejor. Los valores base de la tabla no incluyen la destreza ni objetos mágicos. Dos objetos con nombres distintos pueden dejar la misma CA final.

\begin{reference}{Nombre}{Efecto, uso y límites}
\itemname{leather}{Armadura de cuero} & CA base 8; G, C, L. Una opción inicial para quien no pueda llevar armadura pesada.\\[4pt]
\itemname{scale}{Armadura de escamas} & CA base 6; G, C. La bonificación mágica y la destreza pueden mejorar el resultado.\\[4pt]
\itemname{chain}{Cota de malla} & CA base 5; G, C. Equípala en el cuerpo, no en una mano.\\[4pt]
\itemname{banded}{Armadura de bandas} & CA base 4; G, C. Compara el resultado final en la hoja del personaje.\\[4pt]
\itemname{plate}{Armadura de placas} & CA base 3; G, C. Incluye variantes mágicas o malditas, como la armadura de gran belleza; el nombre no garantiza una mejora.\\[4pt]
\itemname{robe}{Túnica} & Todas las clases. La túnica normal no mejora por sí sola la CA; la Túnica de defensa tiene una bonificación propia.\\[4pt]
\itemname{shield}{Escudo} & Reduce la CA en 1, más su posible bonificación. Debe estar en una mano que pueda usarse; no combina con el uso de un arma a dos manos.\\[4pt]
\itemname{bracers}{Brazales} & Su bonificación puede mejorar la CA al llevarlos en las muñecas. Compara el valor final al cambiar el equipo.\\[4pt]
\itemname{ring-protection}{Anillo de protección} & Llévalo en un dedo. Se utiliza el mejor de los dos anillos y no se suma a una armadura con bonificación mágica.\\[4pt]
\itemname{helmet}{Casco y botas} & Equípalos en las casillas de cabeza y pies. Los cascos y las botas normales no mejoran la CA.\\[4pt]
\end{reference}

\refspace
\subsection*{Herramientas, llaves y objetos especiales}

\begin{reference}{Nombre}{Efecto, uso y límites}
\itemname{spellbook}{Libro de conjuros} & El mago lo necesita para preparar magia. Ponlo en una mano y úsalo con el botón derecho para abrir los conjuros preparados. No se consume al lanzar magia.\\[4pt]
\itemname{holy-symbol}{Símbolo sagrado} & Herramienta del clérigo o paladín. En una mano abre sus oraciones. Llevarlo no sustituye la elección y recuperación de conjuros mediante el descanso.\\[4pt]
\itemname{lockpicks}{Ganzúas} & Coge las ganzúas con el puntero y úsalas sobre una cerradura adecuada. El ladrón puede intentar abrirla; un fallo no significa necesariamente que no sea una cerradura.\\[4pt]
\itemname{key}{Llaves} & Hay llaves de oro, plata, rubí, enjoyadas, de calavera, enanas y drow. Lleva la llave al hueco o cerradura correspondiente. Un color o forma parecidos no garantizan que encaje.\\[4pt]
\itemname{gem}{Gemas} & Conserva las gemas: pueden servir como ofrendas o piezas de un mecanismo. Su uso depende del lugar; no son una moneda que se gaste en una tienda.\\[4pt]
\itemname{stone-object}{Objetos de piedra} & Daga, símbolo sagrado, medallón, collar, orbe, anillo y cetro de piedra. Son piezas de los mecanismos de viaje. Su forma importa; guarda cada una y anota dónde la encontraste.\\[4pt]
\itemname{bones}{Huesos} & Restos humanos, enanos y medianos pueden estar relacionados con personajes de la aventura. Conserva los restos y presta atención a los diálogos; no se usan como comida.\\[4pt]
\itemname{letter}{Carta de encargo y notas} & Pon el documento en una mano y úsalo para leerlo. Las notas dan contexto o pistas. Son documentos, no un automapa que se actualice al caminar.\\[4pt]
\itemname{orb}{Orbe de poder} & Objeto especial de la aventura. No lo trates como una varita recargable ni como un arma ordinaria; conserva los objetos singulares hasta conocer su propósito.\\[4pt]
\itemname{necklace}{Collares y medallones} & Se llevan en el cuello. Algunos son adornos, como el Medallón de la suerte, que no da bonificaciones de combate.\\[4pt]
\end{reference}

\sect{12. Consumibles, varitas y anillos}

\textbf{Consultar un pergamino sin gastarlo.} Cógelo con el botón izquierdo, desde el suelo o desde una casilla del inventario: la zona de mensajes muestra «Pergamino: nombre del conjuro». Funciona con pergaminos de mago y divinos, incluso si el mago ya conoce el conjuro. Devuélvelo a una casilla con otro clic izquierdo. El botón derecho sobre la mano lo utiliza y consume. El nombre permite buscar su efecto en el catálogo de conjuros de este manual.

\textbf{Quién recibe el efecto.} Una poción la bebe quien la sostiene. Un pergamino o una varita que actúan sobre un personaje se aplican al portador al utilizarlos; pásalos antes al destinatario. Los conjuros preparados que piden un destinatario permiten escoger su retrato. No confundas ambos procedimientos.

\textbf{Un uso tiene un coste.} La comida, las pociones y los pergaminos desaparecen al consumirse. Las varitas pierden cargas y desaparecen al agotarse, salvo Silvias, que conserva sus cargas. No se recargan descansando. Un intento que exige una mano libre puede ser rechazado hasta que la liberes.

Al usar un pergamino o una varita, espera a que la mano vuelva a estar disponible.

\textbf{Identificación y duración.} Coge el objeto con el botón izquierdo y lanza Leer magia para identificarlo. Las pociones, varitas y anillos identificados muestran su efecto junto al nombre. Las nubes y los efectos temporales se conservan al guardar y cargar; el descanso hace avanzar su duración. Algunos efectos terminan al descansar, pero los de mayor duración pueden continuar.

\begin{reference}{Nombre}{Efecto, uso y límites}
\itemname{rations}{Raciones / raciones de hierro} & Coge la comida y déjala sobre el plato del inventario del personaje. Se consume y aumenta FOOD hasta un máximo de 100. Hay reservas de distinto tamaño; no todas alimentan lo mismo.\\[4pt]
\itemname{egg}{Huevo de kenku} & También se come sobre el plato. Su aportación de comida varía según el huevo. Conserva provisiones para los descansos largos.\\[4pt]
\itemname{potion-strength}{Poción de fuerza gigante} & Fija la fuerza en 22 durante aproximadamente 150--240 segundos. Después recupera la fuerza anterior, incluso si has bebido otra poción durante el efecto. No cura PV ni veneno.\\[4pt]
\itemname{potion-healing}{Poción de curación} & Restaura 2d4 + 2 PV (4--10), sin superar el máximo. La bebe el personaje que la sostiene.\\[4pt]
\itemname{potion-extra}{Poción de gran curación} & Restaura 3d8 + 3 PV (6--27), sin superar el máximo. No revive a los muertos.\\[4pt]
\itemname{potion-poison}{Poción de veneno} & Puede envenenar al bebedor si no resiste su efecto. Una botella sin identificar no es necesariamente una cura: guarda antes de experimentar.\\[4pt]
\itemname{potion-vitality}{Poción de vitalidad} & Llena FOOD hasta 100. Sacia al personaje, pero no es una poción de curación de PV.\\[4pt]
\itemname{potion-speed}{Poción de velocidad} & Permite usar las armas el doble de rápido durante unos minutos.\\[4pt]
\itemname{potion-invisible}{Poción de invisibilidad} & Aplica invisibilidad al bebedor: los monstruos tienen más dificultad para acertarle. No tiene un plazo fijo; usar una mano o lanzar magia ofensiva puede romperla, igual que el conjuro.\\[4pt]
\itemname{potion-antidote}{Antídoto} & Elimina el envenenamiento. Cura aparte el daño ya sufrido. La poción curativa enana también cura el veneno.\\[4pt]
\itemname{scroll-mage}{Pergamino de mago} & Úsalo desde una mano para gastar su conjuro una vez, o cópialo en el libro mediante Utilizar Pergaminos. Copiarlo también consume el pergamino.\\[4pt]
\itemname{scroll-divine}{Pergamino divino} & Contiene una oración de un solo uso. No se copia en el libro del mago. Si su efecto es sobre un personaje, el uso desde un objeto se aplica a quien lo sostiene.\\[4pt]
\end{reference}

\refspace
\subsection*{Varitas: efectos disponibles}

\begin{reference}{Nombre}{Efecto, uso y límites}
\itemname{wand}{Varita sin poder} & No produce un efecto mágico; aparece el aviso correspondiente.\\[4pt]
\itemname{wand}{Relámpagos} & Lanza Rayo eléctrico. Cada uso válido consume una carga.\\[4pt]
\itemname{wand}{Congelación} & Lanza Cono helado hacia delante. Orienta al grupo antes de usarla.\\[4pt]
\itemname{wand}{Curación} & Usa Curación grave (2d8 + 1 PV) sobre el portador. Pásala al personaje herido antes de activarla.\\[4pt]
\itemname{wand}{Bolas de fuego} & Lanza Bola de fuego. Necesita una trayectoria hacia el enemigo.\\[4pt]
\itemname{wand}{Silvias / Slivias} & Empuja al monstruo de delante una casilla por un paso abierto hasta una casilla vacía. Conserva sus cargas.\\[4pt]
\itemname{wand}{Misiles mágicos} & Lanza Dardo mágico. Es una reserva de ataques cuando faltan conjuros preparados.\\[4pt]
\end{reference}

\refspace
\subsection*{Anillos: llevarlos puestos}

\begin{reference}{Nombre}{Efecto, uso y límites}
\itemname{ring-protection}{Protección} & Su bonificación mejora la CA según las reglas del equipo de la sección 11. Llevarlo en la mochila no basta.\\[4pt]
\itemname{ring}{Sustento} & Equipado en un dedo, evita la pérdida normal de comida, también durante el descanso. No cura heridas.\\[4pt]
\itemname{ring}{Ornamento} & Adorno. No da bonificaciones de combate.\\[4pt]
\itemname{ring}{Hechicería} & Equipado en un dedo, duplica los huecos de preparación de mago para conjuros de nivel 4 y 5. Llevar dos no multiplica de nuevo la bonificación. No concede acceso a niveles superiores ni conocimiento de nuevos hechizos.\\[4pt]
\itemname{ring}{Caída suave} & Equipado en un dedo, evita el daño de caída al que se aplica esta protección. No protege frente a todas las trampas ni frente a los ataques normales.\\[4pt]
\end{reference}

\sect{13. Catálogo de conjuros y oraciones}
El libro identifica los hechizos de mago y el símbolo sagrado, los conjuros divinos: son los dibujos de sus herramientas en el juego, no iconos exclusivos de cada conjuro.

Las tablas distinguen el \textbf{nivel del conjuro} del nivel del personaje. Para preparar una entrada necesitas acceso a ese nivel, un hueco disponible y, si eres mago, conocerla en tu libro. El paladín dispone de una selección propia de oraciones; el menú muestra las que puede preparar.

\textbf{Cómo leer los dados.} 1d8 significa de 1 a 8; 2d8 + 1 significa tirar dos dados de ocho caras y sumar uno (3--17). Son valores base: las defensas, resistencias y tiradas de salvación pueden cambiar el resultado. Una curación nunca supera los PV máximos. Los efectos de protección repetidos no deben darse por acumulativos.

Los personajes multiclase usan el nivel de la clase que lanza el conjuro.

\refspace
\subsection*{Magia de mago}

\refspace
\subsubsection*{Nivel 1}

\begin{reference}{Nombre}{Efecto, uso y límites}
\itemname{spellbook}{Armadura} & Sobre un personaje. Ofrece una CA base de 6, ajustada por destreza, si mejora su protección. Termina cuando el personaje acumula daños por un total de 8 + nivel del lanzador; no absorbe ese daño.\\[4pt]
\itemname{spellbook}{Manos ardientes} & Ataque cercano contra los monstruos de la casilla de delante: 1d3 + 2 por nivel del lanzador. Sitúate frente al enemigo.\\[4pt]
\itemname{spellbook}{Detectar magia} & Mientras dura, los iconos de objetos mágicos se tiñen de azul en el inventario, las manos y el cursor. Al cogerlos también aparece [MAGICO] junto al nombre. No los identifica; Leer magia revela sus propiedades.\\[4pt]
\itemname{spellbook}{Dardo mágico} & Proyectil frontal. Su daño parte de 1d4 + 1 y aumenta con el nivel mediante impactos adicionales. Es práctico desde la retaguardia; las paredes interrumpen la trayectoria.\\[4pt]
\itemname{spellbook}{Leer magia} & Identifica, sin consumirlo, el objeto que llevas en el cursor. Después de lanzarlo se muestra su nombre completo; funciona también con pociones, varitas y anillos. Está disponible para preparar con cualquier mago.\\[4pt]
\itemname{spellbook}{Escudo} & Protección del lanzador. Limita su CA a 4 si era peor y detiene los dardos mágicos que le alcanzan. No es inmunidad a todos los ataques ni a toda la magia.\\[4pt]
\itemname{spellbook}{Golpe eléctrico} & Necesita una mano libre. Crea un ataque mágico en ella, de 1d8 + nivel. Después debes usar esa mano contra un enemigo cercano; no sale un proyectil al elegir el conjuro.\\[4pt]
\end{reference}

\refspace
\subsubsection*{Nivel 2}

\begin{reference}{Nombre}{Efecto, uso y límites}
\itemname{spellbook}{Invisibilidad} & Sobre un personaje. Dificulta que los monstruos lo alcancen, pero no lo hace invulnerable. Usar una mano o lanzar magia ofensiva puede romper la invisibilidad.\\[4pt]
\itemname{spellbook}{Desbloquear} & Abre la puerta ordinaria que tienes delante y libera las puertas que podrían forzarse físicamente. No altera paredes ni cerraduras mágicas o mecanismos de puzle. Está disponible para preparar con cualquier mago que tenga huecos de nivel 2.\\[4pt]
\itemname{spellbook}{Flecha ácida de Melf} & Proyectil frontal: 2d4 por cada tres niveles completos del lanzador, con un mínimo de una tirada. Aplica el daño al impactar, sin daño continuo en turnos posteriores.\\[4pt]
\itemname{spellbook}{Nube ponzoñosa} & Crea una nube en la casilla abierta de delante. Puede inmovilizar brevemente a monstruos y a miembros del grupo que entren en ella. Los no muertos son inmunes. Espera a que se disipe antes de avanzar.\\[4pt]
\end{reference}

\refspace
\subsubsection*{Nivel 3}

\begin{reference}{Nombre}{Efecto, uso y límites}
\itemname{spellbook}{Disolver magia} & Termina los efectos mágicos de un personaje, incluida la invisibilidad y las protecciones beneficiosas. No sustituye a Neutralizar veneno ni a Quitar parálisis.\\[4pt]
\itemname{spellbook}{Bola de fuego} & Proyectil explosivo: 1d6 por nivel del lanzador, antes de las defensas del objetivo. Comprueba la dirección y deja espacio para maniobrar.\\[4pt]
\itemname{spellbook}{Flecha ardiente} & Proyectil frontal de 3d10; pasa a 6d10 al alcanzar el nivel 10 del lanzador. No requiere una flecha física ni un arco.\\[4pt]
\itemname{spellbook}{Velocidad} & Reduce la recuperación total de las armas de unos 3 a 1,5 segundos para los miembros conscientes del grupo durante aproximadamente 30 segundos por (nivel del mago + 3). Si todos los miembros conscientes están acelerados, también se reduce a la mitad la espera de repetición al mantener pulsada una tecla de movimiento.\\[4pt]
\itemname{spellbook}{Paralizar / Bloquear persona} & Proyectil que intenta detener enemigos de tipo humanoide. La resistencia del objetivo puede impedirlo; no sirve contra cualquier criatura.\\[4pt]
\itemname{spellbook}{Invisibilidad 3 m} & Aplica invisibilidad al grupo vivo. Reduce la exposición a los golpes, pero el grupo puede seguir recibiendo daño. Las acciones ofensivas pueden revelar a quien las realiza.\\[4pt]
\itemname{spellbook}{Rayo eléctrico} & Ataque en línea recta, de 1d6 por nivel antes de las defensas. Alinea el pasillo y la orientación del grupo con los enemigos.\\[4pt]
\itemname{spellbook}{Toque vampírico} & Necesita una mano libre y crea un ataque de contacto, con 1d6 por cada dos niveles. Usa después esa mano. Cuando el golpe causa daño, devuelve esa cantidad de PV al atacante; el ataque se gasta al acertar.\\[4pt]
\end{reference}

\refspace
\subsubsection*{Nivel 4}

\begin{reference}{Nombre}{Efecto, uso y límites}
\itemname{spellbook}{Terror} & Intenta hacer huir a los monstruos de la casilla de delante. Puede fallar contra criaturas resistentes. Úsalo para abrir espacio para retirarte.\\[4pt]
\itemname{spellbook}{Tormenta de hielo} & Al alcanzar un objetivo, inflige 3d10 en su casilla y las cuatro adyacentes, formando una cruz. También puede herir al grupo si queda dentro del área. El daño base no depende del nivel del lanzador.\\[4pt]
\itemname{spellbook}{Piel de roca} & Protege a un personaje frente a 1d4 + la mitad del nivel del mago (redondeada hacia abajo) golpes físicos cuerpo a cuerpo que acierten. Cada golpe protegido gasta una carga. No detiene daño de conjuros y puede disiparse.\\[4pt]
\end{reference}

\refspace
\subsubsection*{Nivel 5}

\begin{reference}{Nombre}{Efecto, uso y límites}
\itemname{spellbook}{Nube ácida} & Crea una nube venenosa en la casilla abierta de delante. Puede matar a los enemigos más débiles y herir a los más fuertes. Los no muertos son inmunes. La nube también puede herir a los miembros del grupo que entren en ella.\\[4pt]
\itemname{spellbook}{Cono helado} & Ataca un abanico de casillas frente al grupo. Daño base: 1d4 + 1 por nivel; los enemigos pueden reducirlo con sus defensas. Es útil contra varios adversarios.\\[4pt]
\itemname{spellbook}{Paralizar monstruo} & Versión de inmovilización dirigida a monstruos, no solo humanoides. El proyectil debe llegar al objetivo y su resistencia puede impedir el efecto.\\[4pt]
\end{reference}

\refspace
\subsection*{Magia divina: clérigos y paladines}

\refspace
\subsubsection*{Nivel 1}

\begin{reference}{Nombre}{Efecto, uso y límites}
\itemname{holy-symbol}{Bendición} & Mejora en 1 la tirada para acertar del grupo durante su efecto. No se acumula lanzándola repetidamente sobre un grupo ya bendecido.\\[4pt]
\itemname{holy-symbol}{Curación leve} & Cura 1d8 puntos de vida al personaje elegido, sin superar su máximo. No elimina veneno ni resucita a un muerto.\\[4pt]
\itemname{holy-symbol}{Herida leve} & Ataque de contacto de 1d8. El lanzador tiene que estar en la primera fila y debe haber un enemigo delante; el intento puede fallar.\\[4pt]
\itemname{holy-symbol}{Detectar magia} & Mientras dura, los iconos de objetos mágicos se tiñen de azul en el inventario, las manos y el cursor. Al cogerlos también aparece [MAGICO] junto al nombre. No los identifica; Leer magia revela sus propiedades.\\[4pt]
\itemname{holy-symbol}{Protección contra el mal} & Protege a un personaje: los ataques de monstruos contra él sufren una penalización de 2. Sigue siendo necesario curarlo y evitar quedar rodeado.\\[4pt]
\itemname{holy-symbol}{Imposición de manos} & Facultad del paladín: cura 2 puntos de vida por nivel del lanzador. Se recupera con el descanso; no es una oración normal disponible para todos los clérigos.\\[4pt]
\end{reference}

\refspace
\subsubsection*{Nivel 2}

\begin{reference}{Nombre}{Efecto, uso y límites}
\itemname{holy-symbol}{Ayuda} & Concede 1d8 PV temporales que absorben el daño primero. Los puntos que sobren desaparecen al terminar el conjuro. Volver a lanzar Ayuda mientras está activa no suma más puntos.\\[4pt]
\itemname{holy-symbol}{Hoja llameante} & Necesita una mano libre. Crea un arma temporal de 1d4 + 4, que se utiliza como un arma cuerpo a cuerpo.\\[4pt]
\itemname{holy-symbol}{Paralizar / Bloquear persona} & Intenta inmovilizar humanoides mediante un proyectil. La resistencia puede anularlo. Su uso es equivalente al de la versión de mago.\\[4pt]
\itemname{holy-symbol}{Retrasar veneno} & Sobre un personaje envenenado. Suspende temporalmente el daño del veneno, pero no lo cura. Busca un antídoto antes de que termine.\\[4pt]
\end{reference}

\refspace
\subsubsection*{Nivel 3}

\begin{reference}{Nombre}{Efecto, uso y límites}
\itemname{holy-symbol}{Crear comida} & Llena la reserva de comida de todos los personajes vivos. No restaura puntos de vida ni elimina enfermedades o veneno.\\[4pt]
\itemname{holy-symbol}{Disipar magia} & Termina los efectos mágicos de un personaje, incluida la invisibilidad y las protecciones beneficiosas. No sustituye a Neutralizar veneno ni a Quitar parálisis.\\[4pt]
\itemname{holy-symbol}{Atuendo mágico} & Mejora la clase de armadura a 5, o a un valor menor si quien lo lanza tiene más experiencia. Conserva una protección superior que el personaje ya tuviera.\\[4pt]
\itemname{holy-symbol}{Oración} & Da +1 al ataque del grupo y penaliza en 1 los ataques enemigos. El +1 al ataque no se suma al de Bendición. No cura ni detiene el tiempo.\\[4pt]
\itemname{holy-symbol}{Quitar parálisis} & Libera hasta cuatro personajes paralizados del grupo. No es un antídoto ni una resurrección.\\[4pt]
\end{reference}

\refspace
\subsubsection*{Nivel 4}

\begin{reference}{Nombre}{Efecto, uso y límites}
\itemname{holy-symbol}{Curación grave} & Cura 2d8 + 1 puntos de vida al personaje elegido, hasta su máximo. Reserva las curas grandes para heridas importantes.\\[4pt]
\itemname{holy-symbol}{Herida grave} & Ataque de contacto de 2d8 + 1. Requiere un lanzador en primera fila y un enemigo al alcance.\\[4pt]
\itemname{holy-symbol}{Neutralizar veneno} & Elimina el envenenamiento del personaje elegido. No repone los puntos de vida que ya haya perdido; cúralos por separado.\\[4pt]
\itemname{holy-symbol}{Protección contra el mal 3 m} & Extiende la protección contra ataques de monstruos a los personajes vivos del grupo. Úsala antes de entrar en una zona peligrosa.\\[4pt]
\itemname{holy-symbol}{Protección eléctrica} & No se puede preparar.\\[4pt]
\end{reference}

\refspace
\subsubsection*{Nivel 5}

\begin{reference}{Nombre}{Efecto, uso y límites}
\itemname{holy-symbol}{Curación crítica} & Cura 3d8 + 3 puntos de vida al personaje elegido, sin superar su máximo. No revive a un muerto.\\[4pt]
\itemname{holy-symbol}{Herida crítica} & Ataque de contacto de 3d8 + 3. El lanzador debe estar en la primera fila y necesita un enemigo delante.\\[4pt]
\itemname{holy-symbol}{Columna de fuego} & Ataque mágico de 6d8 al impactar, antes de las defensas del enemigo. Es una opción ofensiva de nivel alto.\\[4pt]
\itemname{holy-symbol}{Resucitar} & Devuelve a un personaje muerto a 1 punto de vida. Los elfos son inmunes. Cura al resucitado antes de volver a combatir.\\[4pt]
\end{reference}

\sect{14. Orientación y pistas}
\textbf{Tu misión, paso a paso.} Explora cada planta, busca una ruta hacia las zonas inferiores, reúne las llaves y piezas necesarias y conserva un grupo capaz de combatir y curarse. El objetivo final es encontrar y derrotar a Xanathar. Las habitaciones secretas y los encargos opcionales pueden recompensar la exploración, pero encontrar todos los secretos no es lo mismo que cumplir la misión principal.

\textbf{Automapa.} La aventura aparece en la salida del V9990 y el mapa en la del MSX. Conecta dos pantallas o cambia de entrada en el monitor. En openMSX, elige la fuente de vídeo MSX para ver el mapa y GFX9000 para volver a la aventura.

El norte queda arriba. La flecha blanca indica dónde está el grupo y hacia dónde mira. Las casillas negras están sin explorar. El mapa marca las puertas visibles en naranja, las escaleras de subida en cian y las de bajada en amarillo. Vuelve a un lugar para ver si algo ha cambiado. Atraviesa las paredes ilusorias para que aparezcan como pasadizos.

El mapa se guarda con el progreso. Abre \texttt{CAMP} para detener la aventura mientras lo consultas.

\textbf{Consultar otras plantas.} \key{F2} muestra la planta anterior, \key{F3} la siguiente y \key{F4} vuelve a la planta en la que estás. También funciona en el campamento. Los huecos abiertos visibles se marcan en rojo y los portales en magenta. Anota por separado las inscripciones y los destinos de los portales.

\subsection*{Primeros minutos}
\begin{enumerate}
\item Coloca a los mejor protegidos delante. Comprueba que llevan sus armas en la mano, no en la mochila.
\item Abre los inventarios, identifica libro, símbolo, comida y munición. Lee la carta de encargo.
\item Prepara curación y algunos ataques mágicos; descansa en un lugar seguro para recuperarlos.
\item Explora un tramo corto, observa brújula y paredes, recoge lo útil y prueba los botones visibles.
\item Guarda antes de bajar por primera vez. Conserva una ruta de regreso a un lugar conocido.
\end{enumerate}

\subsection*{Pistas de menor a mayor ayuda}
\begin{itemize}
\item \textbf{Si parece un callejón sin salida:} mira las paredes laterales, lee las inscripciones y revisa huecos y botones pequeños. No todos los pasos se abren con una llave.
\item \textbf{Si un mecanismo no responde:} piensa en el objeto, su forma y el lugar donde se coloca. Llevar una pieza en la mochila no equivale a usarla en el hueco correcto.
\item \textbf{Si una puerta vuelve a cerrarse:} observa si hay una placa de presión y si necesita peso. Un objeto corriente puede ser más útil en el suelo que en la mochila.
\item \textbf{Si reapareces en un sitio conocido:} puede haber un portal o un cambio de orientación. Comprueba la brújula antes de seguir contando pasos.
\item \textbf{Si encuentras objetos de piedra:} anota su forma y los lugares que tienen huecos semejantes. No los gastes ni abandones sin pensar: ayudan a conectar rutas.
\item \textbf{Si encuentras huesos o personajes amistosos:} escucha antes de atacar y conserva los restos. Algunos encuentros aportan compañeros, ayuda o pistas.
\item \textbf{Si pierdes una pelea:} vuelve con conjuros preparados, una salida libre y curación en las manos adecuadas. Evita rodearte y retírate antes de que caigan los dos de delante.
\end{itemize}

No saltes a un hueco sin preparar una posible caída. Deja reservas en un punto reconocible y anota dónde. Conserva llaves y piezas únicas; deja antes armas corrientes repetidas si falta espacio.

\clearpage
\sect{15. Dudas frecuentes y consulta rápida}
\begin{reference}{Problema}{Comprobación útil}
No puedo atacar & ¿Está el personaje consciente? ¿La mano ha dejado de estar gris? ¿Hay enemigo delante y está en primera fila para un golpe cuerpo a cuerpo?\\
El arco no dispara & Pon flechas en la otra mano y repuestos en el carcaj. Comprueba que el personaje puede usarlo.\\
No puedo ponerme un objeto & Revisa su casilla, las restricciones de clase y si llevas equipo maldito que impida cambiarlo. Una mano ocupada también limita las armas a dos manos.\\
No aparecen conjuros & Comprueba libro o símbolo, personaje, nivel de conjuro, selección en CAMP y descanso completado. Aprendido no significa preparado.\\
La magia pide mano libre & Guarda un objeto de la otra mano. Algunos conjuros crean un ataque que tendrás que usar después desde esa mano.\\
La curación no quita el veneno & Son efectos distintos: antídoto o Neutralizar veneno primero, curación de PV después.\\
No puedo descansar & Busca enemigos cercanos, resuelve el peligro y comprueba comida y estados del grupo.\\
No puedo dejar un miembro & Debe haber al menos cinco personajes y no puedes elegir al primero.\\
Un objeto mágico no hace lo esperado & Consulta su efecto y cómo usarlo en las secciones 11--13.\\
No encuentro mi partida & Comprueba el idioma elegido. Cada idioma tiene su propio guardado. Guarda la partida antes de elegir \texttt{Abandonar Juego}.\\
No se oye la aventura & La fanfarria inicial usa el sonido del MSX; los efectos de aventura y la música de la introducción requieren MoonSound. Revisa Sonidos en Preferencias.\\
\end{reference}

